\documentclass[12pt,a4paper]{article}

\usepackage[utf8]{inputenc}
\usepackage[T1]{fontenc}
\usepackage[brazil]{babel}
\usepackage{geometry}
\usepackage{setspace}
\usepackage{titlesec}
\usepackage{booktabs}
\usepackage{array}
\usepackage{hyperref}
\usepackage{parskip}
\usepackage{microtype}

\geometry{
    top=2.5cm,
    bottom=2.5cm,
    left=3cm,
    right=2.5cm
}

\onehalfspacing

\titleformat{\section}{\large\bfseries}{}{0em}{}[\titlerule]
\titleformat{\subsection}{\normalsize\bfseries}{}{0em}{}

\hypersetup{
    colorlinks=true,
    linkcolor=black,
    urlcolor=black
}

\begin{document}

% ─── Cabeçalho ───────────────────────────────────────────────────────────────
\begin{center}
    {\LARGE\bfseries Monitoramento de Solo via Arduino}\\[0.4em]
    {\large Sistema Distribuído de Sensoriamento para Lavouras de Café}\\[0.8em]
    {\normalsize Benjamin Sergio \quad\textbar\quad \today}
\end{center}

\vspace{0.5em}
\hrule
\vspace{1em}

% ─── Resumo ──────────────────────────────────────────────────────────────────
\begin{abstract}
\noindent
Este documento apresenta o projeto \textit{Monitoramento-Via-Arduino}, um sistema IoT
distribuído para monitoramento de parâmetros de solo em lavouras de café. O sistema
é composto por uma rede de nós sensores de baixo custo e baixo consumo energético
que coletam dados de umidade, temperatura, pH e nutrientes NPK, transmitindo-os
via LoRa a um nó mestre central. O nó mestre agrega os dados, aplica modelos de
inteligência artificial para predição indireta de NPK e os encaminha para plataformas
em nuvem. O protótipo inicial tem orçamento estimado em R\$~250,00.
\end{abstract}

\vspace{0.5em}

% ─── 1. Introdução ───────────────────────────────────────────────────────────
\section{Introdução e Motivação}

O monitoramento preciso das condições do solo é essencial para a gestão eficiente de
lavouras, especialmente em culturas de alto valor como o café. Parâmetros como
umidade, temperatura, pH, condutividade elétrica e concentração de nutrientes
(Nitrogênio, Fósforo e Potássio --- NPK) influenciam diretamente a produtividade e
a qualidade do produto final.

Sistemas comerciais capazes de medir esses parâmetros em campo geralmente apresentam
dois obstáculos: custo elevado e baixa escalabilidade em propriedades extensas.
Soluções baseadas em microcontroladores Arduino e tecnologias de comunicação de longo
alcance (LoRa) surgem como alternativa viável para contornar esses problemas,
possibilitando a implantação de redes distribuídas de sensores a custo acessível.

Este projeto propõe um sistema mestre-escravo no qual múltiplos nós escravos,
distribuídos pela lavoura, coletam dados localmente e os transmitem a um nó mestre
centralizado. O mestre consolida as informações, executa algoritmos de predição
baseados em aprendizado de máquina e as disponibiliza para análise remota via
plataformas em nuvem.

% ─── 2. Arquitetura do Sistema ───────────────────────────────────────────────
\section{Arquitetura do Sistema}

\subsection{Topologia Mestre-Escravo}

O sistema é organizado em dois tipos de nós com responsabilidades distintas:

\textbf{Nó Mestre (ESP32):} Localizado na sede da propriedade, o nó mestre conta
com conectividade Wi-Fi e Bluetooth, processamento central e alimentação robusta.
É responsável por receber os dados dos nós escravos via LoRa, aplicar modelos de
predição de NPK com Random Forest (acurácia $\approx 92\%$) e empacotar os
resultados em JSON para envio à nuvem.

\textbf{Nós Escravos (ATmega328P Pro Mini 3,3~V):} Distribuídos pelo campo,
utilizam o microcontrolador ATmega328P em modo de hibernação profunda para reduzir
o consumo a $0{,}1$--$3\,\mu\text{A}$. Um MOSFET AO3401 realiza o corte físico de
alimentação dos periféricos durante o sono, eliminando os $40$--$80\,\text{mA}$
de corrente parasita que estariam presentes caso o circuito permanecesse energizado.

\subsection{Conjunto de Sensores}

\begin{center}
\begin{tabular}{lll}
\toprule
\textbf{Sensor} & \textbf{Grandezas medidas} & \textbf{Interface} \\
\midrule
JXCT RS485 ModbusJXBS-3001-NPK & Umidade, Temp., CE, NPK & RS485 Modbus \\
BME280 & Temp. ambiente, Umid., Pressão & I2C \\
\bottomrule
\end{tabular}
\end{center}

O sensor RS485 de nível industrial utiliza medição capacitiva indireta para estimar
NPK. Como sensores de NPK diretos de alta precisão têm custo proibitivo, o projeto
adota um modelo de aprendizado de máquina no nó mestre para converter as leituras
brutas de capacitância em valores de NPK com precisão adequada para fins agrícolas.
O BME280 substitui o DHT11/DHT22 por oferecer maior precisão e modo de sono nativo
com consumo inferior.

\subsection{Comunicação LoRa}

O módulo EBYTE E220-900T22D opera na faixa de 433~MHz com potência de até 22~dBm.
A frequência de 433~MHz foi escolhida pela maior penetração em vegetação densa,
característica importante em lavouras de café com dossel fechado. A antena deve ser
posicionada a mais de 2~m acima do dossel para otimizar a propagação do sinal.

A transmissão utiliza estruturas binárias (\texttt{struct}) compactadas, completando
cada ciclo em menos de 5 segundos. O protocolo suporta retransmissão em malha
(\textit{mesh}), permitindo que nós intermediários retransmitam dados de nós mais
distantes até o mestre.

% ─── 3. Prototipagem e Validação ─────────────────────────────────────────────
\section{Prototipagem e Plano de Validação}

\subsection{Protótipo v.01}

O primeiro protótipo tem como objetivo validar a arquitetura de comunicação, medir
o consumo real de energia nos diferentes estados operacionais e verificar a eficácia
do chaveamento por MOSFET. O orçamento estimado é de R\$~250,00, conforme a lista
de componentes abaixo:

\begin{center}
\begin{tabular}{lrr}
\toprule
\textbf{Componente} & \textbf{Qtd.} & \textbf{Custo (BRL)} \\
\midrule
ESP32 DevKit         & 2 & $\approx 50$ \\
Pro Mini 3,3V        & 2 & $\approx 30$ \\
Adaptador FTDI       & 1 & $\approx 15$ \\
BME280               & 3 & $\approx 35$ \\
Ra-02 433~MHz        & 3 & $\approx 40$ \\
MOSFET AO3401 (kit) & 10 & $\approx 20$ \\
INA219 (medidor)     & 1 & $\approx 20$ \\
Protoboard + jumpers & -- & $\approx 40$ \\
\midrule
\textbf{Total}       &   & $\approx 250$ \\
\bottomrule
\end{tabular}
\end{center}

\subsection{Plano de Testes}

Os testes são divididos em quatro eixos:

\textbf{1. Sistema de Comunicação:} Validação da comunicação bidirecional
mestre-escravo e da retransmissão escravo-escravo; análise de perda de pacotes com
e sem elevação da antena; verificação de integridade e autenticação de mensagens.

\textbf{2. Sistema de Energia:} Medição do consumo em cada estado operacional
(espera, ativo, medição, transmissão e hibernação profunda) com o INA219;
avaliação da eficácia do corte por MOSFET em comparação ao simples
modo \textit{sleep} de software.

\textbf{3. Sistema de Medição:} Tempo de aquecimento dos sensores; acurácia e
repetibilidade das leituras; durabilidade em condições de campo (IP68,
exposição a agroquímicos e variações de temperatura).

\textbf{4. Arquitetura Física:} Viabilidade de implantação em campo; otimização do
posicionamento de componentes; resistência ambiental do encapsulamento.

\subsection{Riscos e Mitigações}

Os principais riscos identificados são: (i) atenuação do sinal LoRa por vegetação
densa --- mitigado pela elevação da antena e pelo roteamento em malha; (ii) limitação
da durabilidade dos sensores de solo em campo --- mitigado por encapsulamento IP68 e
protocolos de manutenção preventiva; (iii) imprecisão da predição indireta de NPK
--- mitigado pelo uso de modelos Random Forest treinados com amostras locais e
revalidação periódica.

% ─── Conclusão ───────────────────────────────────────────────────────────────
\section{Considerações Finais}

O projeto \textit{Monitoramento-Via-Arduino} demonstra que é possível construir uma
rede distribuída de monitoramento de solo de baixo custo e longa autonomia,
adequada à realidade das propriedades cafeeiras brasileiras. A combinação de
microcontroladores de ultra-baixo consumo, comunicação LoRa de longo alcance e
modelos de inteligência artificial para predição de NPK oferece uma solução
escalável e tecnicamente robusta.

O protótipo v.01, com orçamento de R\$~250,00, permitirá validar as hipóteses
centrais do projeto antes do desenvolvimento da versão de campo. Os resultados
obtidos na fase de prototipagem guiarão as decisões de componentes, firmware e
encapsulamento para o produto final.

\vspace{1em}
\hrule
\vspace{0.5em}
{\small \textit{Repositório:} \url{https://github.com/benjaminsergio/monitoramento-via-arduino}}

\end{document}
